\chapter{合成器基础}
\label{chap:synth-basics}

\classquote{不徐不疾，得之于手而应于心。}{《庄子·天道》}{不快不慢，手上做得顺、心里也对得上；合成器的参数没有唯一正确值，要在听感与手感之间找平衡。}

本章说明内置复音合成器的信号链、九个可投影参数，以及六个工厂预设的具体取值。
参数名、范围和默认值取自引擎的设备投影代码；数值不做四舍五入以外的加工。

\section{信号链}

内置合成器每轨使用一个定长声部池，最多同时容纳 16 个在响声部。每个声部有两条
振荡器支路、一个四极梯形低通和一个振幅包络。两条支路的波形来自同一张五张波表：
\texttt{hollow}、\texttt{organ}、\texttt{vocal}、\texttt{metallic}、\texttt{glass}。
两条支路各自有电平与失谐，混合后进入低通与包络。

\begin{center}
\begin{tikzpicture}[
  node distance=9mm and 11mm,
  box/.style={draw, rounded corners=2pt, minimum width=24mm, minimum height=8mm, align=center, font=\small},
  sum/.style={draw, circle, minimum size=6mm, inner sep=0pt, font=\small},
  arrow/.style={-{Latex[length=2mm]}, thick}
]
  \node[box] (osc1) {OSC 1\\波表 + 电平 + 失谐};
  \node[box, below=of osc1] (osc2) {OSC 2\\波表 + 电平 + 失谐};
  \node[sum, right=of osc1, yshift=-8mm] (mix) {+};
  \node[box, right=of mix] (filter) {四极梯形\\低通};
  \node[box, right=of filter] (env) {振幅包络\\固定 ADSR};
  \node[box, right=of env] (out) {到插入链\\与总线};

  \draw[arrow] (osc1) -- (mix);
  \draw[arrow] (osc2) -- (mix);
  \draw[arrow] (mix) -- (filter);
  \draw[arrow] (filter) -- (env);
  \draw[arrow] (env) -- (out);
\end{tikzpicture}
\end{center}

滤波器是四极梯形低通。截止频率的有效域由器件钳在 20 Hz 到采样率的 0.45 倍之间；
共振与驱动都是 0 到 1 的旋钮值。驱动为 0 时失真级逐位旁通，为 1 时是 +36 dB
预增益，中间按等 dB 分布；输出侧有自补偿，拧驱动主要换谐波，不换整体音量。

振幅包络是固定的 ADSR：起振 5 ms、衰减 80 ms、 sustain 0.7、释放 50 ms。
当前界面没有提供可调的 ADSR 旋钮，也没有 LFO、调制矩阵或独立包络段控件；
写手册时不应把其他合成器的控件名套到这里。后续若接线，应先在代码里找到对应的
参数投影，再补写本章。

\section{九个可投影参数}

内置合成器从设备链读取的参数名大小写不敏感。下表列出九个参数、有效域与默认值。
“不写”与“写 0”是两件事：不写 \texttt{osc2\_level} 表示第二条支路关，写 0 也是关；
不写 \texttt{cutoff\_hz} 则整个滤波器旁通，写 0 会被器件钳到 20 Hz。

\begin{table}[htbp]
\centering
\small
\begin{tabular}{@{}llll@{}}
\hline
参数名 & 含义 & 有效域 & 默认（不写时） \\
\hline
\texttt{cutoff\_hz}（别名 \texttt{cutoff}） & 低通截止频率 & 20 Hz–0.45×采样率 & 滤波器旁通 \\
\texttt{resonance}（别名 \texttt{res}） & 低通共振 & 0–1 & 0.0 \\
\texttt{drive} & 滤波器驱动/失真 & 0–1 & 0.0 \\
\texttt{osc1\_wave} & 第一条支路波表 & 0–4（四舍五入后取整） & 0（hollow） \\
\texttt{osc2\_wave} & 第二条支路波表 & 0–4（四舍五入后取整） & 0（hollow） \\
\texttt{osc1\_level} & 第一条支路电平 & 0–1 & 1.0（满） \\
\texttt{osc1\_detune\_cents} & 第一条支路失谐 & ±4800 音分 & 0.0 \\
\texttt{osc2\_level} & 第二条支路电平 & 0–1 & 0.0（关） \\
\texttt{osc2\_detune\_cents} & 第二条支路失谐 & ±4800 音分 & 0.0 \\
\hline
\end{tabular}
\caption{内置复音合成器的九个参数}
\label{tab:synth-params}
\end{table}

\screenshot{synth-device-params}{合成器九参数与设备卡读数}

波表下标与名字的对应是固定的：0 = \texttt{hollow}，1 = \texttt{organ}，
2 = \texttt{vocal}，3 = \texttt{metallic}，4 = \texttt{glass}。非有限值或负值按 0
处理；大于 4 的整数由器件钳到最后一张表，不会变成静音。两条支路的电平为 0 时，
该支路一帧也不执行；这是逐位旁通，不是把电平乘成 0。

\section{六个工厂预设的参数}

同一组参数在六个预设里的取值如下。两表合起来就是每个预设的全部九个参数；
表中没有“包络”或“LFO”列，因为当前预设不包含这些参数。

\begin{table}[htbp]
\centering
\small
\begin{tabular}{@{}lrrr@{}}
\hline
预设 & \texttt{cutoff\_hz} & \texttt{resonance} & \texttt{drive} \\
\hline
Hollow Lead & 2400 & 0.20 & 0.00 \\
Warm Organ & 1800 & 0.15 & 0.10 \\
Choir Pad & 2200 & 0.25 & 0.00 \\
FM Bell & 4200 & 0.30 & 0.15 \\
Glass Pluck & 3600 & 0.20 & 0.00 \\
Sub Bass & 420 & 0.35 & 0.25 \\
\hline
\end{tabular}
\caption{工厂预设的滤波器参数}
\label{tab:preset-filter}
\end{table}

\begin{table}[htbp]
\centering
\small
\begin{tabular}{@{}lrrrrrr@{}}
\hline
预设 & \texttt{osc1\_wave} & \texttt{osc2\_wave} & \texttt{osc1\_level} & \texttt{osc1\_detune} & \texttt{osc2\_level} & \texttt{osc2\_detune} \\
\hline
Hollow Lead & 0 & 0 & 0.90 & $-4$ & 0.35 & $+5$ \\
Warm Organ & 1 & 1 & 0.80 & $-5$ & 0.60 & $+6$ \\
Choir Pad & 2 & 0 & 0.85 & $-8$ & 0.55 & $+8$ \\
FM Bell & 3 & 4 & 0.85 & $-3$ & 0.50 & $+11$ \\
Glass Pluck & 4 & 4 & 0.85 & $-4$ & 0.45 & $+4$ \\
Sub Bass & 0 & 1 & 0.95 & $0$ & 0.35 & $-5$ \\
\hline
\end{tabular}
\caption{工厂预设的振荡器参数；失谐单位为音分}
\label{tab:preset-osc}
\end{table}

从表 \ref{tab:preset-filter} 与表 \ref{tab:preset-osc} 可以直接读出每个预设的
性格：Hollow Lead 是一条九成电平的 hollow 加一条低电平同波表，适合柔和主音；
Warm Organ 让两条 organ 支路同时开，靠失谐做出风琴的厚度；Choir Pad 用 vocal
作主支路、hollow 作辅支路；它的铺底听感来自 vocal 波表与较低截止的组合；
FM Bell 与 Glass Pluck 都用 metallic/glass 波表的高截止，前者更亮、后者更短促；
Sub Bass 把截止压到 420 Hz，只留低频。

\subsection{三台标志合成器：减法合成的三种性格}

参数表能告诉你“现在是什么”，却不能告诉你“为什么这样弹”。减法合成的历史里有三台琴，正好对应本章的三组动作：Minimoog 把“少”做成重量，Juno 把“薄”做成宽度，TB-303 把“一条线”做成律动。

\textbf{Minimoog（1970）}把原本模块化、占房间的 Moog 合成器塞进一台可携带的琴：三个振荡器、梯形低通、没有记忆预设。它留下的不是“参数越多越好”，而是“每条支路都要有角色”。本章只有两条支路，正好练这件事：先让 \texttt{osc1} 站住，再用 \texttt{osc2} 的电平与失谐补厚度；第三层留给编曲，不要硬塞进一个音色。夜半的 Sub Bass 用 420 Hz 低通加一条 organ 支路，就是“少”的另一种写法。

\textbf{Juno-60/Juno-106（1982/1984）}用 DCO 与合唱，把“薄”做成“宽”。它的合唱不是真的加一条第三振荡器，而是用调制把时间与左右错开；本章没有合唱效果，只能在两条同波表支路上做小失谐，例如 $\pm5$ 到 $\pm9$ 音分，再把宽度留给后续效果。不要声称这等于 Juno 的合唱，它是“用现有参数做近似”，不是复刻。

\textbf{TB-303（1981）}本来是为贝斯手做的单振荡器加滤波器，却在舞池里被当成 acid 的主音。它提醒我们：共振与截止不是音色装饰，它们会改变音高感与律动。把 \texttt{resonance} 从 0.2 推到 0.5 时，截止附近的谐波被强调，旋律的重心会跟着移动；这不是音量，别用 \texttt{drive} 去补。夜半没有 303 的滑音与 accent 控件，能练的是“一条线＋一个滤波器”的判断。

三台琴的共通点：先把一条线做清楚，再加第二条；先听相位与失谐，再谈加厚。参数是手，听感才是耳朵；不要让预设名替你听。

\section{本章与帮助主题的映射}

\begin{table}[htbp]
\centering
\small
\begin{tabular}{@{}p{0.34\linewidth}p{0.52\linewidth}@{}}
\hline
本章小节 & 应用内帮助锚点 \\
\hline
信号链 & \texttt{synth-signal-chain} \\
九个可投影参数 & \texttt{synth-parameters} \\
六个工厂预设的参数 & \texttt{synth-factory-presets} \\
\hline
\end{tabular}
\caption{第二章帮助主题映射}
\end{table}

\section*{延伸阅读}
\begin{itemize}[leftmargin=1.5em]
\item \url{https://www.soundonsound.com/techniques}——Sound On Sound 技术栏目：用行业教程对照振荡器、滤波器与包络的实际听感。
\item \url{https://www.ableton.com/en/blog/}——Ableton 博客：查看合成器音色拆解与制作人工作流，补足参数背后的音乐判断。
\end{itemize}
